% !TEX root = ../my-thesis.tex

\chapter{Fazit und Ausblick}
\label{sec:conclusion}

\section{Fazit}
\label{sec:conclusion:fazit}

Diese Arbeit hat fünf Workflow-Management-Systeme anhand zweier
Workflow-Muster und derselben Rechenlogik verglichen, zunächst
unter kontrollierten Bedingungen auf einem einzelnen Rechner und
anschließend auf dem HPC-System MOGON. Untersucht wurden die
Modellierung, das Laufzeitverhalten anhand von Makespan, Overhead
und Koordinationsgrößen sowie der Betrieb unter Slurm.

Beide Muster ließen sich in allen fünf Systemen umsetzen. Die
wesentlichen Unterschiede liegen daher nicht in der Abbildbarkeit,
sondern in der Art der Workflow-Beschreibung, im
Koordinationsverhalten und im Betriebsmodell. Beim Overhead
bildeten StreamFlow, Nextflow und Merlin auf beiden Ebenen eine
Gruppe mit niedrigen Werten, AiiDA lag darüber und Pegasus in der
lokalen Untersuchung mit großem Abstand an der Spitze. Da der
Overhead nicht mit der Rechenlast wächst, verliert er mit
zunehmender Aufgabengröße an Gewicht.

Auf dem HPC-System zeigte sich ein starker Einfluss der
Ausführungsumgebung auf die gemessenen Laufzeiten. Die vom Cluster verursachten Zeitanteile können die Unterschiede zwischen den Systemen vollständig überlagern, und die Systeme binden den Ressourcenverwalter
unterschiedlich ein, von einzelnen Jobs je Schritt bis zur
gemeinsamen Worker-Allokation. Aussagekräftige Vergleiche
erforderten deshalb die Ausführung innerhalb einer bestehenden
Allokation.

Für die Eignungsfrage folgt daraus, dass sie sich nicht an der
Abbildbarkeit der Muster festmachen lässt, denn diese war in allen
fünf Systemen gegeben. Die Unterschiede lagen im Laufzeitverhalten
und im Betrieb. Beim Laufzeitverhalten zeigte sich zugleich, dass
der Overhead mit wachsender Rechenlast an Gewicht verliert, sodass
die gemessenen Unterschiede bei hinreichend großen Rechenschritten
kaum noch ins Gewicht fallen. Ab diesem Punkt entscheiden andere Eigenschaften über die Wahl, etwa die Ausrichtung auf hybride Ausführungsumgebungen bei
StreamFlow, die Abarbeitung vieler gleichartiger Aufgaben innerhalb
einer Allokation bei Merlin oder die dauerhafte Nachvollziehbarkeit
bei AiiDA. Ein in allen Belangen überlegenes System ergab sich nicht,
die Wahl richtet sich nach der Struktur des Workflows, der Größe
der Aufgaben und den Anforderungen des Anwendungsfalls.

\section{Ausblick}
\label{sec:conclusion:ausblick}

An den Grenzen der Untersuchung setzen mögliche weiterführende
Arbeiten an. Der Vergleich ließe sich zunächst auf komplexere
Workflow-Strukturen ausdehnen und mit Rechenlasten wiederholen, die
realen wissenschaftlichen Anwendungen näherkommen, etwa mit großen
Datenmengen und intensiven Dateizugriffen. So ließe sich prüfen, ob
die hier beobachteten Unterschiede zwischen den Systemen auch unter
solchen Bedingungen fortbestehen.

Ein zweiter Ansatzpunkt ist die Skalierung. Die Untersuchung
umfasste höchstens vier parallele Verarbeitungsinstanzen. Mit
deutlich mehr Aufgaben ließe sich prüfen, wie sich die Systeme bei
großen Aufgabenzahlen verhalten, insbesondere Merlin, dessen Entwurf
genau darauf zielt. Messungen auf weiteren HPC-Systemen würden
zudem zeigen, welche der auf MOGON beobachteten Betriebs- und
Laufzeiteigenschaften an die konkrete Ausführungsumgebung gebunden
sind.

Offen bleibt schließlich der HPC-Vergleich für Pegasus. Lässt sich
Pegasus mit der benötigten HTCondor-Umgebung auf einem Cluster
betreiben, kann die quantitative Untersuchung auf alle fünf Systeme
ausgeweitet werden.